\section{Related Work}

\subsection{Respiratory-Audio Datasets and Benchmark Models}

Coswara contains multiple respiratory and vocal tasks with participant
metadata. COUGHVID is a large, heterogeneous cough corpus
\cite{bhattacharya2023coswara,orlandic2021coughvid}. The COVID-19 Sounds resource
similarly includes breathing, cough, and voice and introduced realistic
participant-independent baselines \cite{xia2021covid}. These resources differ
in recruitment, label construction, class prevalence, sound prompts, and
quality control. Dataset identity is therefore part of the experimental
condition, not an interchangeable filename.

Early systems used MFCCs, spectral descriptors, OpenSMILE functionals, support
vector machines, convolutional networks, and score fusion. Feature-ranking
work showed that the useful representation can depend on the dataset
\cite{meister2022ranking}, and the DiCOVA study of Avila \emph{et al.} combined
6,373-dimensional ComParE features with a log-mel CNN
\cite{avila2021feature}. Multimodal Coswara work reported gains from averaging
predictions across respiratory and vocal tasks, but selected combinations from
many possible subsets \cite{chetupalli2023multimodal}. These studies motivate multimodal
modeling while also making validation-only fusion selection essential.

\subsection{Deep and Transformer Representations}

Transfer learning with CNN, LSTM, and ResNet representations has been evaluated
for cough, breathing, and speech \cite{pahar2022transfer}. WavLM provides a
self-supervised transformer representation learned from large speech corpora
\cite{chen2022wavlm}. Aytekin \emph{et al.} proposed the hierarchical
spectrogram transformer (HST), which partitions a $224\times224$ log-mel
spectrogram into patches and expands the attended context through local-window
attention and patch merging \cite{aytekin2024hst}. Their reported internal
results are important architecture references, but their Cambridge and
COUGHVID experiments use dataset-specific tasks and 10-fold partitions. They
are not direct estimates of Coswara-to-COUGHVID transfer.

The ESWA DNDT/DNDF study combines 193 cough features, recursive feature
elimination with cross-validation, Bayesian hyperparameter optimization,
SMOTE, and threshold moving \cite{islam2026robust}. It reports high
within-dataset AUCs for Coswara and COUGHVID, but its explicit cross-dataset
table also contains much lower source-to-target values. This distinction is
central to the present comparison: within-dataset cross-validation and frozen
external transfer answer different questions.

\subsection{Bias, Drift, and Clinical Utility}

Han \emph{et al.} demonstrated that participant overlap and demographic
imbalance can produce optimistic audio estimates \cite{han2022realistic}.
Coppock \emph{et al.} extended this analysis through prespecified matched and
longitudinal cohorts, symptom baselines, incremental models, and utility
analysis \cite{coppock2024audio}. Ganitidis \emph{et al.} studied model drift
and adaptation across development and later cough data \cite{ganitidis2025drift}.
Together, these studies shift the research question from maximum internal AUC
to transportability under a stated target population and time.

Prediction-model reporting guidance reaches the same conclusion. TRIPOD+AI
requires explicit data provenance, participant flow, modeling choices,
performance measures, and limitations \cite{collins2024tripod}; PROBAST+AI
emphasizes risk of bias and applicability across participants, predictors,
outcomes, and analysis \cite{moons2025probast}. We therefore report AUROC with
AUPRC, calibration, threshold-dependent metrics, and uncertainty, and keep
confirmatory comparisons separate from exploratory mechanism analyses.
